
\subsubsection{6.1 Interpreting Near-Perfect Coupling}

Near-perfect correlations (r > 0.99, p < $2\times 10^{-17})$ across TrustfulQA, AdvBench, and BOLD challenge the foundational assumption that trustworthiness dimensions measure orthogonal properties. Results suggest two non-mutually-exclusive explanations:

\textbf{Unified Capability Hypothesis:} All three benchmarks primarily measure general model quality---the same underlying factor driving overall performance. Models with poor training data diversity, weak calibration, or brittle representations fail across all dimensions proportionally. This aligns with prior observations that RLHF improves multiple dimensions simultaneously without dimension-specific targeting. However, correlation cannot prove causal unity---r > 0.99 could reflect shared confounds (e.g., all benchmarks sensitive to training data diversity) rather than unified constructs. \textbf{Discriminating prediction:} 10-benchmark expansion shows r > 0.95 for >80\% of pairs AND factor analysis reveals single component explaining >90\% variance.

\textbf{Insufficient Resolution Hypothesis:} Distinct failure modes exist but require >3 benchmarks to distinguish. Our k = 2 clustering (TrustfulQA+AdvBench versus BOLD) may reflect measurement artifact rather than conceptual structure. HELM's 50+ benchmark coverage suggests broader evaluation could resolve dimensions invisible in our 3-benchmark subset. \textbf{Discriminating prediction:} 10-benchmark expansion shows mean r < 0.85 with $\geq 30$\% of pairs < 0.7 AND clustering produces well-separated groups (silhouette > 0.5).

The two hypotheses make opposing predictions testable through broader measurement. Current data (r > 0.99 from 3 benchmarks) fits both hypotheses---distinguishing them requires expansion to 5--10 benchmarks and factor analysis.

\subsubsection{6.2 Bootstrap-Silhouette Divergence: Methodological Lesson}

Clustering analysis revealed perfect stability (100\% bootstrap consistency) but poor separation (silhouette = 0.274), a divergence that illuminates the difference between two cluster quality metrics:
\begin{itemize}
\item \textbf{Bootstrap consistency} measures reproducibility: Do benchmarks cluster together across resampling? (100\% = yes)
\item \textbf{Silhouette score} measures discriminant validity: Are clusters far apart in distance space? (0.274 = no)
\end{itemize}

These properties are orthogonal. With r > 0.99 correlations, all benchmarks lie nearly collinear in correlation space, yielding tiny inter-benchmark distances (0.003--0.007 range). Ward linkage can partition this space into stable groups (100\% consistency), but silhouette---which compares within-cluster cohesion to between-cluster separation---penalizes small absolute distances regardless of stability.

\textbf{Implication for benchmark design:} Clustering quality thresholds (silhouette > 0.5) may be unattainable for high-correlation data (r > 0.95), even when structure is stable. Alternative metrics (Calinski-Harabasz, Davies-Bouldin) or revised thresholds (silhouette > 0.2 for r > 0.95 data) may better suit trustworthiness evaluation.

\subsubsection{6.3 Limitations}

\textbf{L1: Sample Size (3 Benchmarks) --- Critical:} Hierarchical clustering with k $\geq$ 3 requires n $\geq$ 3 samples, limiting analysis to k = 2 evaluation. Cannot validate the full taxonomy claim. Expanding to 5--10 benchmarks (ToxiGen, BBQ, HellaSwag, GSM8K, HumanEval) would enable k = 3--5 clustering with adequate sample size.

\textbf{L2: Extremely High Correlations (r > 0.99) --- High:} Near-perfect correlations produce minimal distance variation (0.003--0.007 range), making silhouette > 0.5 thresholds unattainable even for stable clusters. Reporting both silhouette (separation) and bootstrap consistency (stability) distinguishes reproducibility from discriminant validity. Future work should test alternative metrics (Calinski-Harabasz, Davies-Bouldin) or revise thresholds (silhouette > 0.2 for r > 0.95 data).

\textbf{L3: Cluster Interpretability --- Medium:} k = 2 solution (TrustfulQA+AdvBench versus BOLD) lacks clear semantic interpretation. No obvious mapping to epistemic uncertainty, distribution shift, or bias amplification mechanisms. Larger benchmark suite may yield semantically coherent clusters (e.g., reasoning benchmarks versus safety benchmarks versus fairness benchmarks).

\textbf{L4: Data Source Noise --- Acknowledged:} Public leaderboard aggregation may introduce measurement error from inconsistent evaluation protocols. However, r > 0.99 correlations suggest signal dominates noise.

\subsubsection{6.4 Implications}

\textbf{For Model Developers:} If r > 0.99 coupling generalizes beyond the 3-benchmark sample, interventions targeting one dimension (e.g., calibration training for reliability) plausibly improve others (robustness, fairness) simultaneously. This suggests multi-dimensional co-training may be more efficient than dimension-specific fixes, though intervention validation remains future work.

\textbf{For Benchmark Designers:} Results reveal that 3-benchmark evaluation cannot resolve failure mode taxonomy despite perfect cluster stability. Future trustworthiness benchmarks should aim for 5--10 dimensions minimum to enable robust clustering (k = 3--5 evaluation) with adequate sample size.

\textbf{For Evaluation Frameworks (HELM, BIG-bench):} Aggregation-based reporting should be complemented by correlation analysis. Discovering that three benchmarks correlate at r > 0.99 informs resource allocation---if benchmarks measure near-redundant constructs, evaluation budgets could shift toward broader coverage (new dimensions) rather than deeper coverage (more tasks per dimension).
